\paragrafos{Γραφική παράσταση της αναλογίας: μια ημιευθεία από την αρχή}{Α.6.4, «Γραφική παράσταση σχέσης αναλογίας»}

\begin{anagke}
Ο Νίκος και ο Βασίλης κάνουν ποδήλατο στον παραλιακό δρόμο. Ο Νίκος κάνει $18$ χιλιόμετρα την ώρα και ο Βασίλης $12$. Η Πολυξένη γράφει τις αποστάσεις τους σε πίνακες: ώρα $0,\ 1,\ 2,\ 3$ και χιλιόμετρα $0,\ 18,\ 36,\ 54$ για τον Νίκο, και $0,\ 12,\ 24,\ 36$ για τον Βασίλη. Ο Ηλίας κοιτάζει τους αριθμούς και λέει: «Βλέπω ότι ο Νίκος είναι γρηγορότερος, αλλά δεν μπορώ να δω \emph{πόσο}».

Η Αγγελική παίρνει ένα τετραγωνισμένο φύλλο και, όπως στη ναυμαχία, σημειώνει τα ζευγάρια ώρα--χιλιόμετρα ως σημεία. Όλα τα σημεία του Νίκου πέφτουν σε μια ευθεία γραμμή που ξεκινά από την αρχή. Το ίδιο και του Βασίλη, αλλά η δική του γραμμή είναι πιο «ξαπλωτή». «Όσο πιο γρήγορα, τόσο πιο απότομη η γραμμή», παρατηρεί η Ρήνη. Ο Στέλιος ρωτάει αν αυτό ισχύει για κάθε αναλογία, και η Φανή δοκιμάζει με το μέλι του παππού και μια δική της σχέση, που δεν πέφτει σε ευθεία.

Ένα σχήμα λέει με μια ματιά αυτό που ένας πίνακας λέει με πολλά νούμερα. Σήμερα βρίσκουμε τι είδους σχήμα έχει η αναλογία και πώς το χρησιμοποιούμε για να διαβάζουμε τιμές.
\end{anagke}

\begin{lexi}{γραφική παράσταση, ημιευθεία}{γραφική παράσταση}
Η \lexiword{γραφική παράσταση} είναι η «εικόνα που γράφεται» με σημεία και γραμμές. Η \lexiword{ημιευθεία} είναι η μισή ευθεία: ξεκινά από ένα σημείο και συνεχίζει στο άπειρο προς μία κατεύθυνση. Στις αναλογίες το σημείο εκκίνησης είναι η αρχή $O$ των αξόνων.
\end{lexi}

\begin{metafrasi}
Ένα ζεύγος τιμών $(x,\,y)$ ενός πίνακα. & ένα σημείο του επιπέδου\\
Ανάλογα ποσά με συντελεστή $\alpha$. & σημεία πάνω στην ημιευθεία $y=\alpha x$ από το $O$\\
Μεγαλύτερος συντελεστής. & πιο απότομη ημιευθεία\\
Για $x=1$ η τεταγμένη. & είναι ο συντελεστής $\alpha$\\
\end{metafrasi}

\tmima{Η θεωρία}

Αν δύο ποσά είναι ανάλογα με συντελεστή $\alpha$, δηλαδή $y=\alpha x$, ένα σημείο του επιπέδου αντιστοιχεί σε κάθε ζεύγος τιμών $(x,\,y)$. Για παράδειγμα με $\alpha=1{,}5$ έχουμε τον πίνακα
\begin{center}
\renewcommand{\arraystretch}{1.3}
\begin{tabular}{c|cccccc}
$x$ & $0$ & $1$ & $2$ & $4$ & $5$ & $6$\\\hline
$y=1{,}5x$ & $0$ & $1{,}5$ & $3$ & $6$ & $7{,}5$ & $9$
\end{tabular}
\end{center}
και τα σημεία $(0,\,0)$, $(1,\,1{,}5)$, $(2,\,3)$, $(4,\,6)$, $(5,\,7{,}5)$, $(6,\,9)$.

\begin{center}
\begin{tikzpicture}[x=0.9cm,y=0.5cm]
\draw[gray!35,very thin] (0,0) grid (6.4,9.4);
\draw[thick,-{Stealth}] (0,0)--(6.8,0) node[below] {$x$};
\draw[thick,-{Stealth}] (0,0)--(0,10) node[left] {$y$};
\foreach \i in {1,...,6}{\node[below,font=\scriptsize] at (\i,0) {\i};}
\foreach \j in {1,3,5,7,9}{\node[left,font=\scriptsize] at (0,\j) {\j};}
\draw[gblue,thick] (0,0)--(6,9);
\foreach \x/\y in {0/0,1/1.5,2/3,4/6,5/7.5,6/9}{\fill[gblue] (\x,\y) circle (2.3pt);}
\node[gblue,right] at (6,9) {$y=1{,}5x$};
\end{tikzpicture}
\end{center}

Παρατηρούμε ότι τα σημεία βρίσκονται πάνω στην ίδια ημιευθεία με αρχή το $O(0,\,0)$.

\begin{idiotita}{γραφική παράσταση της αναλογίας}
Τα σημεία που αντιστοιχούν στα ζεύγη τιμών $(x,\,y)$ δύο \textbf{ανάλογων} ποσών βρίσκονται πάνω σε μία \textbf{ημιευθεία που ξεκινά από την αρχή $O(0,\,0)$ των ημιαξόνων}. Και αντίστροφα, αν τα σημεία ενός πίνακα βρίσκονται σε ημιευθεία από το $O$, τα ποσά είναι ανάλογα.
\end{idiotita}

Ο συντελεστής $\alpha$ βρίσκεται από τη γραφική παράσταση: είναι η τεταγμένη του σημείου που έχει τετμημένη $1$, ή γενικά το πηλίκο $\frac{y}{x}$ οποιουδήποτε σημείου της. Όσο μεγαλύτερος ο $\alpha$, τόσο πιο «απότομη» η ημιευθεία. Από τη γραφική παράσταση διαβάζουμε και τιμές που δεν υπάρχουν στον πίνακα: ανεβαίνουμε από ένα $x$ ως την ημιευθεία και διαβάζουμε το $y$.

\tmimasmall{Πώς διαπιστώνω γραφικά ότι τα ποσά \emph{δεν} είναι ανάλογα}

Αν τα σημεία δεν είναι σε ευθεία (έχουν «τσακίσματα» ή κάμπτονται), ή αν είναι σε ευθεία που δεν περνά από το $O$, τα ποσά δεν είναι ανάλογα. Π.χ. η $y=x+2$ έχει σημεία $(0,\,2)$, $(1,\,3)$, $(2,\,4)$ που βρίσκονται σε ευθεία που κόβει τον $Oy$ στο $2$.

\begin{prosoxi}
\begin{itemize}[leftmargin=1.2em,itemsep=0.3ex,topsep=2pt]
\item Δεν αρκεί να είναι τα σημεία σε ευθεία: πρέπει η ευθεία να περνά από την αρχή $O$. Η $y=x+2$ είναι ευθεία χωρίς να είναι αναλογία.
\item Δεν ενώνουμε τα σημεία «όπως μας βολεύει»: ενώνουμε με ημιευθεία μόνο αν ξέρουμε ότι υπάρχουν και ενδιάμεσες τιμές (π.χ. μπορεί να πάρω $2{,}5$ κιλά μέλι).
\item Προσοχή στη μονάδα στους άξονες: κάθε διαίρεση πρέπει να είναι ίση παντού, αλλά δεν χρειάζεται να είναι ίδια στον $x$ και στον $y$.
\item Σε ανάλογα ποσά, η μεγάλη τιμή του $y$ δεν σημαίνει αυτόματα μεγάλο συντελεστή: μετράει ο λόγος $\frac{y}{x}$.
\end{itemize}
\end{prosoxi}

\tmima{Λυμένα παραδείγματα}

\paradeigma{Σχεδιάζω την αναλογία}
\tixero Τα $x,\,y$ είναι ανάλογα με $y=1{,}5x$.\par
\tithelo Να τα παραστήσω γραφικά.\par
\possindeo Κάνω πίνακα ($x=0,\,2,\,4,\,6$ και $y=0,\,3,\,6,\,9$), τοποθετώ τα σημεία και φέρνω ημιευθεία από το $O$ που τα περιέχει.\par
\apantisi Η ημιευθεία από το $O$ που περνά από το $(2,\,3)$ (σχήμα θεωρίας).\par
\elegxos Το $(4,\,6)$ είναι πάνω της και $\frac{6}{4}=1{,}5$.

\paradeigma{Ποιος πίνακας είναι αναλογία;}
\tixero Πίνακες με $x=0,\,1,\,2,\,3$: (Α) $y=0,\,2,\,3,\,5$. (Β) $y=2,\,4,\,6,\,8$. (Γ) $y=0,\,2,\,4,\,6$. (Δ) $y=0,\,1,\,4,\,9$.\par
\tithelo Να διαπιστώσω ποιοι δίνουν ανάλογα ποσά.\par
\possindeo Σε κάθε πίνακα κοιτάζω αν τα σημεία είναι σε ημιευθεία από το $O$. Α: οι λόγοι $\frac{2}{1}=2$, $\frac{3}{2}=1{,}5$: όχι. Β: έχει $y=2$ όταν $x=0$, άρα δεν ξεκινά από το $O$: όχι. Γ: $\frac{y}{x}=2$ παντού: ναι. Δ: $\frac{1}{1}=1$, $\frac{4}{2}=2$: όχι, είναι καμπύλη.\par
\apantisi Μόνο ο Γ ($y=2x$).\par
\elegxos Το Β είναι ευθεία $y=x\cdot2+2$ που περνά από το $(0,\,2)$.

\paradeigma{Οι δύο ποδηλάτες}
\tixero Ο Νίκος κάνει $18$ km/h και ο Βασίλης $12$ km/h.\par
\tithelo Γραφική παράσταση για κάθε ποδηλάτη, απόσταση στις $2$ ώρες, και πότε ο Νίκος κάνει $54$ km.\par
\possindeo Οι αποστάσεις είναι $y=18x$ και $y=12x$ ($x$ ώρες, $y$ χλμ.). Τα σημεία: Νίκος $(0,\,0)$, $(1,\,18)$, $(2,\,36)$, $(3,\,54)$. Βασίλης $(0,\,0)$, $(1,\,12)$, $(2,\,24)$, $(3,\,36)$. Στις $2$ ώρες: $36$ και $24$ km. Για τα $54$ km ο Νίκος: $54:18=3$ ώρες.\par
\apantisi Η ημιευθεία του Νίκου είναι πιο απότομη. Στις $2$ ώρες: $36$ και $24$ km. $54$ km σε $3$ ώρες. Μετά από $3$ ώρες η διαφορά τους είναι $54-36=18$ km.\par
\elegxos Ο συντελεστής $18>12$ κάνει την ημιευθεία του Νίκου πιο απότομη.

\begin{center}
\begin{tikzpicture}[x=1.5cm,y=0.06cm]
\draw[gray!35,very thin] (0,0) grid[xstep=0.5,ystep=10] (3.2,60);
\draw[thick,-{Stealth}] (0,0)--(3.5,0) node[below] {ώρες};
\draw[thick,-{Stealth}] (0,0)--(0,66) node[left] {km};
\foreach \i in {1,2,3}{\node[below,font=\scriptsize] at (\i,0) {\i};}
\foreach \j in {20,40,60}{\node[left,font=\scriptsize] at (0,\j) {\j};}
\draw[gblue,thick] (0,0)--(3,54) node[right] {Νίκος};
\draw[gblue!60!black,thick,dashed] (0,0)--(3,36) node[right] {Βασίλης};
\end{tikzpicture}
\end{center}

\paradeigma{Βρίσκω τον συντελεστή από το γράφημα}
\tixero Δύο ανάλογα ποσά έχουν γραφική παράσταση ημιευθεία από το $O$ που περνά από το σημείο $(4,\,10)$.\par
\tithelo Ο συντελεστής και το $y$ όταν $x=7$.\par
\possindeo $\alpha=\frac{10}{4}=2{,}5$. Τότε $y=2{,}5\cdot7=17{,}5$.\par
\apantisi $\alpha=2{,}5$ και $y=17{,}5$.\par
\elegxos Το σημείο $(1,\,2{,}5)$ είναι κι αυτό στην ημιευθεία.

\tmima{Ασκήσεις}

\zoni{Το κατάλαβα}{σημεία, ημιευθείες, πίνακες}

\ex Φτιάξε πίνακα για την $y=2x$ με $x=0,\,1,\,2,\,3,\,4,\,5$ και σχεδίασε τη γραφική παράσταση.

\ex Σχεδίασε στο ίδιο σύστημα αξόνων τις $y=x$, $y=2x$ και $y=0{,}5x$ για $x$ από $0$ ως $6$. Ποια είναι πιο απότομη και γιατί;

\ex Ποιοι από τους πίνακες παριστάνουν ανάλογα ποσά; α) $x$: $0,2,4,6$, $y$: $0,3,6,9$ \quad β) $x$: $0,1,2,3$, $y$: $1,2,3,4$ \quad γ) $x$: $0,1,2,3$, $y$: $0,1,3,6$ \quad δ) $x$: $0,2,4,5$, $y$: $0,5,10,12{,}5$.

\ex Τα σημεία $(2,\,3)$, $(4,\,6)$, $(6,\,9)$ και $(8,\,12)$ βρίσκονται σε ημιευθεία από το $O$; Βρες τον συντελεστή. Και τα $(2,\,3)$, $(4,\,6)$, $(5,\,8)$;

\ex Η ημιευθεία από το $O$ περνά από το $(5,\,12)$. Βρες τον συντελεστή αναλογίας, το $y$ για $x=2$ και το $x$ για $y=30$.

\ex Για την ευθεία $y=x+2$ φτιάξε πίνακα για $x=0,\,1,\,2,\,3$ και σχεδίασε. Γιατί δεν είναι αναλογία;

\zoni{Το χρησιμοποιώ}{μέλι, ποδήλατα, βενζίνη}

\ex \textbf{Το μέλι.} Σχεδίασε τη γραφική παράσταση της τιμής για μέλι $8$ €/κιλό. Διάβασε από αυτή το κόστος για $2{,}5$ κιλά και τα κιλά που παίρνεις με $36$ €.

\ex \textbf{Οι ποδηλάτες.} Ο Νίκος κάνει $18$ km/h και ο Βασίλης $12$ km/h. Σχεδίασε τις αποστάσεις για $0$ ως $4$ ώρες. Πόση είναι η διαφορά τους μετά από $3$ ώρες; Σε πόση ώρα ο Βασίλης κάνει $30$ km;

\ex \textbf{Η βενζίνη.} Σχεδίασε την τιμή της βενζίνης προς $1{,}85$ €/λίτρο για $0$ ως $40$ λίτρα (χρησιμοποίησε κατάλληλη μονάδα στον κάθε άξονα). Πόσο κοστίζουν $12$ λίτρα περίπου;

\ex \textbf{Το σχοινί.} Ένα σχοινί βάρους $0{,}4$ kg ανά μέτρο. Σχεδίασε το βάρος για $0$ ως $5$ μ. και πες πόσο ζυγίζουν $3{,}5$ μ.

\zoni{Το σκέφτομαι}{γιατί ημιευθεία από το $O$}

\ex Γιατί η γραφική παράσταση μιας αναλογίας περνά πάντα από το $O$;

\ex Η γραφική παράσταση ενός πίνακα είναι ευθεία που δεν περνά από το $O$. Μπορούν τα ποσά να είναι ανάλογα; Εξήγησε.

\ex Δύο ανάλογα ποσά έχουν σημεία $(3,\,4)$ και $(6,\,8)$. Ποιο είναι το $y$ για $x=9$ και ποιο το $x$ για $y=2$;

\ex Ο Στέλιος ένωσε τα σημεία της θερμοκρασίας του πυρετού της Φανής με ημιευθεία από το $O$. Γιατί είναι λάθος;

\zoni{Για το σπίτι}{μέτρησε και σχεδίασε}

\ex Μέτρησε το μήκος $1,\,2,\,3,\,4$ βημάτων σου και σχεδίασε το γράφημα βήματα--μέτρα. Είναι αναλογία;

\ex Πάρε τα κιλά και την τιμή τριών προϊόντων από το σούπερ μάρκετ και σχεδίασε την τιμή σε σχέση με το βάρος για ένα από αυτά.
